%------------------------------------------------------------------------------%
\section{Séries de Taylor}

Nessa seção definiremos as Séries de Taylor e mostraremos como calcular seus coeficientes.
Um ponto importante é que construir uma série de Taylor para uma função não garante
a convergência da série para essa função, a série pode, por exemplo,
convergir para outra função. Essa questão será discutida na próxima seção.
Apresentamos primeiro a definição dos Polinômios de Taylor que podem ser
construídos para quaisquer funções que possuam o número necessário
de derivadas no ponto $a$.

%------------------------------------------------------------------------------%
\begin{definicao}{Polinômio de Taylor}{def:polinomio-taylor}
	Se uma função real $f$ tem derivadas até ordem $n$ em uma vizinhança de $a$,
	definimos o
	\emph{Polinômio de Taylor}\index{Polinômio!Taylor}\index{Taylor!polinômio}
  de $f$ e ordem $n$ centrado em $a$ como
	\begin{align*}
	  P_n(x) & = \sum_{k=0}^n \frac{\,f^{(k)}(a)\,}{k!}(x-a)^k \\[3mm]
	         & = f(a)
             + f'(a)(x-a)
             + \frac{\,f''(a)\,}{2!}(x-a)^2
             + \cdots
             + \frac{\,f^{(n)}(a)\,}{n!}(x-a)^n     \;.
	\end{align*}
\end{definicao}
%------------------------------------------------------------------------------%

Quando a função possuir derivadas de todas as ordens em $a$ podemos definir sua Série de Taylor.

%------------------------------------------------------------------------------%
\begin{definicao}{Série de Taylor}{def:serie-taylor}
  Dada uma função real $f$ infinitamente derivável em uma vizinhança de $a$,
  sua \emph{Série de Taylor}\index{Série!Taylor}\index{Taylor!série}
  centrada em $a$, é dada por
  \begin{multline*}
    \sum_{n=0}^\infty \frac{\,f^{(n)}(a)\,}{n!}(x-a)^n = \\
     f(a) +
     f'(a)(x-a) +
     \frac{\,f''(a)\,}{2!}(x-a)^2 +
     \cdots +
     \frac{\,f^{(n)}(a)\,}{n!}(x-a)^n +
     \cdots
  \end{multline*}
\end{definicao}
%------------------------------------------------------------------------------%

A Série de Maclaurin é um caso particular da Série de Taylor.

%------------------------------------------------------------------------------%
\begin{definicao}{Série de Maclaurin}{def:serie-maclaurin}
  A \emph{Série de Maclaurin}\index{Série!Maclaurin}\index{Maclaurin!série}
é uma Série de Taylor centrada na origem, isso é, $a=0$
\[
	\sum_{n=0}^\infty \frac{f^{(n)}(0)}{n!}x^n
    = f(0)
    +	f'(0)\,x
    +	\frac{\,f''(0)\,}{2!}x^2
    + \cdots
    +	\frac{\,f^{(n)}(0)\,}{n!}x^n
    + \cdots
\]
\end{definicao}
%------------------------------------------------------------------------------%

Observe que estamos usando a notação de derivada zero $f^{(0)}$
para indicar a própria função $f$, isso é,
\[
  f^{(0)}(x) = f(x)
\]
Apresentamos como primeiro exemplo o cálculo da Série de Taylor da
função $e^x$ centrada na origem, isso é, sua Série de Maclaurin.

%------------------------------------------------------------------------------%
\begin{example}[ex:serie-taylor-exponencial]
	Calcule a Série de Maclaurin da função $f(x)=e^x$.
	\solution
	O primeiro passo é calcular as derivadas da função.
	Esse exemplo foi escolhido pela simplicidade em calcular as derivadas
  da função exponencial que têm sempre a mesma expressão
	\[
		f^{(k)}(x) = \frac{d^ke^x}{dx^k} = e^x     \;.
	\]
	Como queremos a série centrada em $x=0$, precisamos avaliar as derivadas nesse ponto
	\[
		f^{(k)}(0) = e^0 = 1     \;.
	\]
	Podemos agora calcular os coeficientes da Série de Taylor
	\begin{align*}
	  c_0 & = f(0) = 1                            \;,\\[3mm]
	  c_1 & = f'(0) = 1                           \;,\\[3mm]
	  c_2 & = \frac{f^{(2)}(0)}{2!} = \frac{1}{2} \;,\\[3mm]
	  c_3 & = \frac{f^{(3)}(0)}{3!} = \frac{1}{6} \;.
	\end{align*}
	Repetindo essas operações observamos que
	\[
		c_k = \frac{f^{(k)}(0)}{k!} = \frac{1}{k!} \;.
	\]
	Podemos agora construir a Série de Taylor da função exponencial
	\begin{align*}
		S(x) & = \sum_{n=0}^\infty \frac{f^{(n)}(a)}{n!}(x-a)^n \\[3mm]
         & = \sum_{n=0}^\infty \frac{f^{(n)}(0)}{n!}x^n     \\[3mm]
         & = \sum_{n=0}^\infty \frac{x^n}{n!}
           = 1 + x + \frac{x^2}{2} + \frac{x^3}{3!} + \cdots + \frac{x^n}{n!}+\cdots
	\end{align*}
\end{example}
%------------------------------------------------------------------------------%

\begin{figure}
  \begin{center}%
    \vspace{-0.2\baselineskip}%
    \scalebox{0.9}{% Title: gl2ps_renderer figure
% Creator: GL2PS 1.4.2, (C) 1999-2020 C. Geuzaine
% For: Octave
% CreationDate: Tue Sep 28 14:08:54 2021
\setlength{\unitlength}{1pt}
\begin{picture}(0,0)
\includegraphics[width=17cm,height=7cm]{fig_taylor_exp}
\end{picture}%
\begin{picture}(481,198)(0,0)
\fontsize{10}{0}\selectfont\put(21.9528,28.661){\makebox(0,0)[t]{\textcolor[rgb]{0.15,0.15,0.15}{{-3}}}}
\fontsize{10}{0}\selectfont\put(107.622,28.661){\makebox(0,0)[t]{\textcolor[rgb]{0.15,0.15,0.15}{{-2}}}}
\fontsize{10}{0}\selectfont\put(193.291,28.661){\makebox(0,0)[t]{\textcolor[rgb]{0.15,0.15,0.15}{{-1}}}}
\fontsize{10}{0}\selectfont\put(364.63,28.661){\makebox(0,0)[t]{\textcolor[rgb]{0.15,0.15,0.15}{{1}}}}
\fontsize{10}{0}\selectfont\put(450.299,28.661){\makebox(0,0)[t]{\textcolor[rgb]{0.15,0.15,0.15}{{2}}}}
\fontsize{10}{0}\selectfont\put(283.956,74.207){\makebox(0,0)[l]{\textcolor[rgb]{0.15,0.15,0.15}{{1}}}}
\fontsize{10}{0}\selectfont\put(283.956,112.272){\makebox(0,0)[l]{\textcolor[rgb]{0.15,0.15,0.15}{{2}}}}
\fontsize{10}{0}\selectfont\put(283.956,150.337){\makebox(0,0)[l]{\textcolor[rgb]{0.15,0.15,0.15}{{3}}}}
\fontsize{10}{0}\selectfont\put(283.956,188.403){\makebox(0,0)[l]{\textcolor[rgb]{0.15,0.15,0.15}{{4}}}}
\fontsize{10}{0}\selectfont\put(119.616,9.49606){\makebox(0,0)[l]{\textcolor[rgb]{0,0,0}{{$P_1$}}}}
\fontsize{10}{0}\selectfont\put(66.5008,9.49606){\makebox(0,0)[l]{\textcolor[rgb]{0,0,0}{{$P_3$}}}}
\fontsize{10}{0}\selectfont\put(3.9622,9.49606){\makebox(0,0)[l]{\textcolor[rgb]{0,0,0}{{$P_5$}}}}
\fontsize{10}{0}\selectfont\put(11.6725,98.9494){\makebox(0,0)[r]{\textcolor[rgb]{0,0,0}{{$P_2$}}}}
\fontsize{10}{0}\selectfont\put(11.6725,140.821){\makebox(0,0)[r]{\textcolor[rgb]{0,0,0}{{$P_4$}}}}
\end{picture}
}%
    \vspace{-\baselineskip}%
  \end{center}
  \caption{Comparação da função exponencial com seus polinômios de Taylor.}
  \label{fig:funcao_tayplor_exponencial}
\end{figure}

Ainda não podemos dizer que a série calculada nesse exemplo seja igual a função~$e^x$.
Entretanto ela aproxima muito bem a função como podemos ver na
Figura~\ref{fig:funcao_tayplor_exponencial},
onde a função $e^x$ está mostrada em azul. Nesse gráfico também estão as primeiras
somas parciais da série, correspondentes aos Polinômio de Taylor com graus de 1 a 5.
O gráfico em magenta exibe $P_5$, observe que a partir de $x=-1$ esse gráfico
sobrepõe o da função exponencial.

\OnlineBox{https://colab.research.google.com/drive/1X3XFXqi-ZF6KbKYTA0oJUYREw9EnhlL_?usp=sharing}
{
  Código Python para ilustrar algumas séries de Taylor.
}

%------------------------------------------------------------------------------%
\subsection*{Séries Binomiais}

Vamos agora apresentar um caso particular de grande importância histórica e
prática a Série Binomial.
Começamos apresentando o
\emph{Binômio de Newton}\index{Binômio de Newton}\index{Newton!binômio}
que é a expansão de um binômio elevado a uma potência inteira, $(a-b)^m$,
estamos interessados no caso particular
\begin{equation}
  \label{eq:binomio-newton}
  (1+x)^m = \sum_{k=0}^m \binom{m}{k} \, x^k
  \qquad\qquad \forall\; m\in \N
\end{equation}
onde
\begin{equation}
  \label{eq:combinacoes}
  \binom{m}{k}
  = \frac{m!}{\,k!(m-k)!\,}
  = \frac{m(m-1)\cdots(m-k+1)}{k!}
\end{equation}
é o \emph{Número de Combinações}\index{Número de combinações}
de $m$ elementos tomados $k$ a $k$. Observe que, para qualquer~$m$
\[
  \binom{m}{0} = 1    \;.
\]
Podemos verificar a fórmula~\eqref{eq:binomio-newton}
calculando a Serie de Taylor da função
\[
  f(x) = (1+x)^m
\]
onde $m$ é um número inteiro positivo.
Começamos calculando as derivadas de $f$ em $x=0$
\begin{align*}
  f(x)   & = (1+x)^m           & f(0)   & = 1      \\[2mm]
  f'(x)  & = m(1+x)^{m-1}      & f'(0)  & = m      \\[2mm]
  f''(x) & = m(m-1)(1+x)^{m-2} & f''(0) & = m(m-1) \;.
\end{align*}
Repetindo o processo observamos que para todo $k$
\[
  f^{(k)}(x) = m(m-1)\cdots (m-k+1)(1+x)^{m-k}
\]
portanto
\[
  f^{(k)}(0) = m(m-1)\cdots (m-k+1)  \;.
\]
Note que para todo $k>m$ teremos que $f^{(k)}(0)=0$.
Dessa forma, a Série de Taylor de $f$ centrada em zero é
uma soma finita
\[
  \sum_{k=0}^\infty \frac{\;f^{(k)}(0)\;}          {k!}\,x^k =
  \sum_{k=1}^m      \frac{\;m(m-1)\cdots (m-k+1)\;}{k!}\,x^k
\]
onde os coeficientes são exatamente o
número de combinações de $m$ elementos tomados $k$ a $k$.
Assim, podemos escrever que
\begin{equation}
  \label{eq:binomio-newton-2}
  \sum_{k=0}^\infty \frac{\;f^{(k)}(0)\;}{k!}\,x^k =
  \sum_{k=0}^m \binom{m}{k} x^k
  \qquad\qquad m \in \N    \;.
\end{equation}
Como o somatório é finito temos apenas uma igualdade entre polinômios
e podemos escrever que a soma é igual a função
\[
  (1+x)^m = \sum_{k=0}^m \binom{m}{k} x^k
  \qquad\qquad m \in \N     \;.
\]

Queremos agora generalizar esse resultado para $m\in \R$.
Essa generalização é a \emph{Série Binomial}\index{Série!binomial}.
Com $m$ real as derivadas de $f$ continuam com as mesmas expressões
\begin{align*}
  f^{(k)}(x) & = m(m-1)\cdots (m-k+1)(1+x)^{m-k} \\
  f^{(k)}(0) & = m(m-1)\cdots (m-k+1)     \;.
\end{align*}
Porém, $f^{(k)}(0)$ deixa de ser zero para $k>m$. Isso transforma a
soma~\eqref{eq:binomio-newton-2} em uma série propriamente dita
\[
  \sum_{k=0}^\infty \frac{\;f^{(k)}(0)\;}{k!}\,x^k =
  \sum_{k=0}^\infty \binom{m}{k} x^k
\]
onde
\[
  \binom{m}{k} = \frac{m(m-1)\cdots(m-k+1)}{k!}
  \qquad\qquad m \in \R     \;.
\]
Observe que como $m$ não é um inteiro positivo não podemos calcular seu
fatorial.
Outro ponto fundamental é que com infinitos termos na Série de Taylor
não podemos afirmar que ela é igual a função sem verificar sua
convergência.
O próximo teorema conclui a generalização da
equação~\eqref{eq:binomio-newton} para todo $m$ real,
mostrando que a Série de Taylor de $f(x)=(1+x)^m$ converge para a função.

%------------------------------------------------------------------------------%
\begin{teorema}{Série Binomial}{teo:serie-binomial}
  Para qualquer número $m\in\R$ e todo $x\in(-1, 1)$, vale que
  \[
    (1+x)^m = \sum_{k=0}^\infty \binom{m}{k} \, x^k
  \]
  onde
  \[
    \binom{m}{k} = \frac{m(m-1)\cdots(m-k+1)}{k!}    \;.
  \]
\end{teorema}
%------------------------------------------------------------------------------%
\begin{proof}
  Vamos aplicar o Teste da Razão nos termos da série
  \[
    a_k = \binom{m}{k} x^k = \frac{m(m-1)\cdots(m-k+1)}{k!} x^k  \;.
  \]
  A razão entre os módulos dos termos é
  \begin{align*}
    \abs{ \frac{a_{k+1}}{a_k} }
    & = \abs{  \binom{m}{k+1} x^{k+1}                    \;
        \left[ \binom{m}{k}   x^k \right]^{-1} }         \\[5mm]
    & = \abs{ \frac    {\,m(m-1)\cdots(m-k)  \,}{(k+1)!} \;
              \frac{k!}{\,m(m-1)\cdots(m-k+1)\,}         \;
              \frac{\,x^{k+1}\,}{x^k} }                  \\[3mm]
    & = \frac{\,\abs{m-k}\,}{k+1} \abs{x}  \;.
  \end{align*}
  O limite dessa razão, quando $k\to\infty$, é igual a $\abs{x}$.
  Então, a série converge absolutamente quando $\abs{x}<1$,
  ou seja, o intervalo de convergência é $(-1, 1)$.
  Isso significa que nesse intervalo podemos definir a função
  \[
    h(x) = \sum_{k=0}^\infty \binom{m}{k} x^k  \;.
  \]
  Nosso objetivo agora é mostrar que
  \[
    h(x) = (1+x)^m
  \]
  Para verificar essa igualdade, vamos usar uma expressão equivalente
  \[
    g(x) = \frac{h(x)}{(1+x)^{m}} = (1+x)^{-m}h(x) = 1
  \]
  Avaliando $g$ em zero temos
  \[
    g(0)
    = (1+0)^{-m}h(0)
    = \sum_{k=0}^\infty \binom{m}{k} 0^k
    = \binom{m}{0} 0^0
    = 1   \;.
  \]
  Agora basta provar que $g$ é uma função constante,
  ou seja, vamos mostrar que $g'(x)=0$ para todo $x\in (-1,\,1)$.
  Calculando a derivada e igualando a zero temos
  \[
    g'(x) = -m(1+x)^{-m-1} h(x) + (1+x)^{-m} h'(x) = 0  \;.
  \]
  Rearranjando a equação
  \begin{align*}
    -m(1+x)^{-m-1} h(x) + (1+x)^{-m} h'(x) & = 0                  \\
    (1+x)^{-m} h'(x)                       & = m(1+x)^{-m-1} h(x) \\
    (1+x)^{-m} (1+x)^{m+1} h'(x)           & = m h(x)             \\
    (1+x) h'(x)                            & = m h(x)   \;.
  \end{align*}
  obtemos que $g'(x)$ será zero se
  \[
    (1+x) h'(x) = m h(x)
  \]
  Para demonstrar essa igualdade vamos partir da expressão do lado esquerdo
  e mostrar que ela é igual a do lado direito.
  Começamos por derivar a série $h(x)$ termo a termo
  \[
    h'(x) = \sum_{k=1}^\infty \binom{m}{k} k x^{k-1}   \;.
  \]
  Podemos agora escrever
  \begin{align*}
    (1+x) h'(x) & = h'(x) + xh'(x) \\
                & = \sum_{k=1}^\infty \binom{m}{k} k x^{k-1}
                + \sum_{k=1}^\infty \binom{m}{k} k x^k    \;.
  \end{align*}
  Para somar as duas séries precisamos igualar o
  expoente de $x$, para isso alteramos o índice da primeira série
  \[
    (1+x) h'(x)
       = \sum_{k=0}^\infty \binom{m}{k+1} (k+1) x^k
       + \sum_{k=1}^\infty \binom{m}{k}    k    x^k \\[3mm]
  \]
  Agora igualamos os índices das séries separando o termo $k=0$ da primeira
  e agrupamos os somatórios
  \[
    (1+x) h'(x) = m + \sum_{k=1}^\infty
                      \bigg[ \binom{m}{k+1} (k+1)
                           + \binom{m}{k}    k    \;
                      \bigg] x^k     \;.
  \]
  Vamos agora substituir as fórmulas para as combinações
  \begin{align*}
    (1+x) h'(x)
    & = m + \sum_{k=1}^\infty \bigg[\; \frac{m(m-1)\cdots(m-k)  }{(k+1)!} \;(k+1)             \\[3mm]
    &       \hspace{23mm}            + \frac{m(m-1)\cdots(m-k+1)}{ k!   } \; k \;\bigg] x^{k} \;.
  \end{align*}
  Podemos agora simplificar os coeficientes de $x^k$ e colocar em evidência a parte comum
  \begin{align*}
    (1+x) h'(x)
    & = m + \sum_{k=1}^\infty \bigg[\; \frac{m(m-1)\cdots(m-k)  }{k!    }                \\[3mm]
    &       \hspace{23mm}            + \frac{m(m-1)\cdots(m-k+1)}{(k-1)!} \;\bigg] x^{k} \\[3mm]
    & = m + \sum_{k=1}^\infty \frac{m(m-1)\cdots(m-k+1)}{(k-1)!}
                              \left[ \frac{(m-k)}{k}+1 \right] x^{k}                     \\[3mm]
    & = m+\sum_{k=1}^\infty \frac{m(m-1)\cdots(m-k+1)}{(k-1)!}\;\frac{\,m\,}{k}\; x^{k}  \\[3mm]
    & = m+\sum_{k=1}^\infty \frac{m(m-1)\cdots(m-k+1)}{k!} \;m x^{k}          \;.
  \end{align*}
  Colocando a constante $m$ em evidencia e notando que a
  fração é a expressão da combinação temos que
  \begin{align*}
    (1+x) h'(x)
      & = m + m \sum_{k=1}^\infty \binom{m}{k} x^{k}                             \\[3mm]
      & = m \left[ 1            + \sum_{k=1}^\infty \binom{m}{k} x^{k} \right]   \\[3mm]
      & = m \left[ \binom{m}{0} + \sum_{k=1}^\infty \binom{m}{k} x^{k} \right]   \\[3mm]
      & = m \sum_{k=0}^\infty \binom{m}{k} x^{k}                                 \\[3mm]
      & = m h(x)          \;.
  \end{align*}
  Isso mostra que, para $x\in(-1, 1)$,
  \[
    (1+x)^m = \sum_{k=0}^\infty \binom{m}{k} x^k          \;.
  \]
\end{proof}
%------------------------------------------------------------------------------%

As \emph{Funções Potência}\index{Função!potência},
\(
  f(x) = x^r
\),
são muito comuns em diversas áreas e aplicações da Matemática.
Se o expoente $r$ for um número inteiro elas também são fáceis para avaliar.
Porém, em casos como
\[
f(x) = x^{\sqrt{2}}     \qquad\qquad
f(x) = \sqrt[3]{x}      \qquad\qquad
f(x) = x^{\sfrac{5}{7}}
\]
precisamos aproximar o valor da função por seu Polinômio de Taylor.
O próximo exemplo mostra como calcular a Série de Taylor para essas
funções.

%------------------------------------------------------------------------------%
\begin{example}
  Encontre a Série de Taylor da função $\sqrt[3]{x}$.
  \solution
  Primeiro escrevemos a função na forma binomial
  \[
    \sqrt[3]{x}
    = x^{\sfrac{1}{3}}
    = \big(1 + (x-1)\big)^{\sfrac{1}{3}}
    = (1+y)^{\sfrac{1}{3}}
  \]
  onde $y=x-1$.
  Usando o Teorema~\ref{teo:serie-binomial} podemos escrever
  \[
    (1+y)^{\sfrac{1}{3}} = \sum_{k=0}^\infty \binom{\sfrac{1}{3}}{k}\, y^k
    \qquad y \in(-1, 1)          \;.
  \]
  Desfazendo a mudança de variáveis, concluímos que, para $x \in(0, 2)$,
  \begin{align*}
     \sqrt[3]{x}
     & = \sum_{k=0}^\infty \binom{\sfrac{1}{3}}{k} (x-1)^k \\[3mm]
     & = \!\binom{\sfrac{1}{3}}{0} (x-1)^0
       + \!\binom{\sfrac{1}{3}}{1} (x-1)^1
       + \!\binom{\sfrac{1}{3}}{2} (x-1)^2
       + \!\binom{\sfrac{1}{3}}{3} (x-1)^3
       + \cdots \\[3mm]
     & = 1 + \frac{1}{3}(x-1) \\[3mm]
     &\hphantom{=1\;} + \frac{\sfrac{1}{3}\left(\sfrac{1}{3}-1\right)}{2}(x-1)^2 \\[3mm]
     &\hphantom{=1\;} + \frac{\sfrac{1}{3}\left(\sfrac{1}{3}-1\right)\left(\sfrac{1}{3}-2\right)}{3!}(x-1)^3
       + \cdots \\[3mm]
     & = 1 + \frac{1}{3}(x-1) - \frac{1}{9}(x-1)^2 + \frac{5}{81}(x-1)^3 + \cdots
  \end{align*}
\end{example}
%------------------------------------------------------------------------------%

%------------------------------------------------------------------------------%
